% ============================================================
% Section 3: Greek Mathematics (Slides 15-24)
% ~600 BCE -- 300 CE
% ============================================================
\section{Greek Mathematics}

% ----------------------------------------------------------
% Slide 15: Section Divider
% ----------------------------------------------------------
{
\setbeamercolor{background canvas}{bg=greekBlue}
\begin{frame}[plain]
  \centering
  \vspace{1cm}
  \textcolor{white}{\Huge\bfseries The Greek Revolution}\\[0.8em]
  \textcolor{white!80}{\large Math Becomes Logic\quad $\sim$600 BCE -- 300 CE}\\[1.5em]

  % Mini-map: Greece, Alexandria, Hellenistic world
  \visualplaceholder[5cm]{Mediterranean outline hint}

  \vspace{0.5em}
  % Timeline marker
  \visualplaceholder[5cm]{Diagram}

  \note{
    Section transition. The Greeks did something no civilization before them had done:
    they asked ``WHY does this work?'' and demanded logical proof.
    Point out the geographic spread: Miletus, Samos, Athens, Syracuse, Alexandria.
  }
\end{frame}
}

% ----------------------------------------------------------
% Slide 16: Thales of Miletus
% ----------------------------------------------------------
\begin{frame}{Thales of Miletus}

  \begin{columns}[T]
    \begin{column}{0.55\textwidth}
      \begin{personbox}[Thales of Miletus (Thal\={e}s ho Mil\={e}sios)]{greekBlue}
        \small
        \textbf{Dates:} c.\ 624 -- c.\ 546 BCE\\
        \textbf{Origin:} Miletus, Ionia (modern Turkey)\\
        \textbf{Contribution:} Traditionally regarded as the first person to use
        \emph{deductive reasoning} in mathematics
      \end{personbox}

      \vspace{0.5em}
      \begin{itemize}
        \item \alert{Thales' theorem:} any angle inscribed in a semicircle is a right angle
        \item Credited with predicting a solar eclipse
        \item Marks the shift from ``\emph{what}'' to ``\emph{why}''
      \end{itemize}
    \end{column}

    \begin{column}{0.42\textwidth}
      % Thales' theorem diagram
      \visualplaceholder[5cm]{Semicircle}

      \vspace{0.5em}
      \visualplaceholder[4.5cm]{Map showing Miletus on the coast of Ionia (modern Turkey)}
    \end{column}
  \end{columns}

  \note{
    Thales is traditionally the ``first mathematician'' in the Greek sense ---
    the first to demand logical justification.
    \verifytag{The eclipse prediction story may be legendary. Exact dates
    are uncertain --- c.\ 624--546 BCE are traditional.}
    ``Before Thales, math told you WHAT. After Thales, math also told you WHY.''
  }
\end{frame}

% ----------------------------------------------------------
% Slide 17: Pythagoras and the Pythagoreans
% ----------------------------------------------------------
\begin{frame}{Pythagoras and the Pythagoreans}

  \begin{columns}[T]
    \begin{column}{0.52\textwidth}
      \begin{personbox}[Pythagoras of Samos (Pythagoras ho Samios)]{greekBlue}
        \small
        \textbf{Dates:} c.\ 570 -- c.\ 495 BCE\\
        \textbf{Origin:} Samos, Greece; active in Croton, Magna Graecia (southern Italy)\\
        \textbf{Contribution:} Founded a school that saw numbers as the
        foundation of reality
      \end{personbox}

      \vspace{0.3em}
      \begin{itemize}
        \item $a^2 + b^2 = c^2$ bears his name (but known earlier --- recall Plimpton 322)
        \item Discovery that $\sqrt{2}$ is \alert{irrational} --- a crisis for their philosophy
        \item ``All is number'' --- numbers govern the universe
      \end{itemize}
    \end{column}

    \begin{column}{0.45\textwidth}
      % Visual proof of Pythagorean theorem
      \visualplaceholder[5cm]{Right triangle}

      \vspace{0.3em}
      {\small\centering $16 + 9 = 25$ \quad\checkmark}
    \end{column}
  \end{columns}

  \note{
    ``When they discovered irrational numbers, it reportedly caused a philosophical
    crisis --- legend says they tried to keep it secret.''
    \verifytag{Whether Pythagoras himself or his followers made these discoveries
    is debated; ancient sources attribute much to ``the Pythagoreans'' collectively.}
    Remind students: we saw Pythagorean triples on Plimpton 322, 1,200 years earlier.
  }
\end{frame}

% ----------------------------------------------------------
% Slide 18: Euclid -- The Elements
% ----------------------------------------------------------
\begin{frame}{Euclid --- \emph{The Elements}}

  \begin{columns}[T]
    \begin{column}{0.55\textwidth}
      \begin{personbox}[Euclid of Alexandria (Eukleides)]{greekBlue}
        \small
        \textbf{Dates:} fl.\ c.\ 300 BCE (birth/death unknown)\\
        \textbf{Origin:} Active in Alexandria, Ptolemaic Egypt\\
        \textbf{Contribution:} Wrote \emph{The Elements} --- the most influential
        mathematics textbook in history
      \end{personbox}

      \vspace{0.4em}
      \begin{itemize}
        \item Used for \textbf{over 2,000 years} as a textbook
        \item Organized geometry into a \alert{deductive system from axioms}
        \item Proved the \textbf{infinitude of prime numbers}
        \item 13 books covering geometry, number theory, proportion
      \end{itemize}
    \end{column}

    \begin{column}{0.42\textwidth}
      \visualplaceholder[4.5cm]{Page from a medieval manuscript of \emph{The Elements} (Bodleian Library or Vatican Library)}

      \vspace{0.3em}
      % Euclid's proof structure
      \visualplaceholder[5cm]{Diagram}
    \end{column}
  \end{columns}

  \note{
    Euclid's \emph{Elements} is the second most printed book in Western history
    (after the Bible). Abraham Lincoln taught himself logic by reading it.
    The axiomatic method Euclid used is STILL the foundation of how we do mathematics today.
    ``Every proof you write in a geometry class follows Euclid's template.''
  }
\end{frame}

% ----------------------------------------------------------
% Slide 19: Archimedes -- The Greatest?
% ----------------------------------------------------------
\begin{frame}{Archimedes --- The Greatest?}

  \begin{columns}[T]
    \begin{column}{0.52\textwidth}
      \begin{personbox}[Archimedes of Syracuse]{greekBlue}
        \small
        \textbf{Dates:} c.\ 287 -- 212 BCE\\
        \textbf{Origin:} Syracuse, Sicily (Greek colony)\\
        \textbf{Contribution:} Calculated $\pi$ with high accuracy; discovered
        principles of buoyancy and the lever; method of exhaustion
      \end{personbox}

      \vspace{0.3em}
      \begin{itemize}
        \item Bounded $\pi$: \  $3\tfrac{10}{71} < \pi < 3\tfrac{1}{7}$
        \item Used inscribed / circumscribed polygons
        \item \textbf{Method of exhaustion} \textrightarrow\ compared to early integral calculus
        \item Killed during the siege of Syracuse (212 BCE)
      \end{itemize}
    \end{column}

    \begin{column}{0.45\textwidth}
      % Archimedes polygon method
      \visualplaceholder[5cm]{Circumscribed hexagon}

      \vspace{0.3em}
      {\small\centering
        $\underbrace{P_{\text{in}}}_{\text{lower}} < \underbrace{C}_{\pi} < \underbrace{P_{\text{out}}}_{\text{upper}}$\\[0.3em]
        Used 96-sided polygons!
      }

      \vspace{0.3em}
      \visualplaceholder[4cm]{The Archimedes Palimpsest (Walters Art Museum, public domain)}
    \end{column}
  \end{columns}

  \note{
    ``Archimedes was killed by a Roman soldier during the siege of Syracuse,
    reportedly saying `Do not disturb my circles.'\,''
    \verifytag{This is from Plutarch and may be legendary.}
    \verifytag{Characterizing the method of exhaustion as ``anticipating integral calculus''
    is debated. Some historians see it as a genuine precursor; others argue it is
    fundamentally different in concept. Present as a comparison, not an identity.}
  }
\end{frame}

% ----------------------------------------------------------
% Slide 20: Eratosthenes -- Measuring the Earth
% ----------------------------------------------------------
\begin{frame}{Eratosthenes --- Measuring the Earth}

  \begin{columns}[T]
    \begin{column}{0.48\textwidth}
      \begin{personbox}[Eratosthenes of Cyrene]{greekBlue}
        \small
        \textbf{Dates:} c.\ 276 -- c.\ 194 BCE\\
        \textbf{Origin:} Cyrene (modern Libya); worked in Alexandria\\
        \textbf{Contribution:} Calculated Earth's circumference; Sieve of Eratosthenes for primes
      \end{personbox}

      \vspace{0.4em}
      \begin{itemize}
        \item Compared shadow angles at \textbf{Alexandria} and \textbf{Syene} (Aswan)
        \item At summer solstice: no shadow in Syene, $\sim 7.2^{\circ}$ shadow in Alexandria
        \item Result: $\sim$\textbf{252,000 stadia} $\approx$ remarkably close to the true value
      \end{itemize}
    \end{column}

    \begin{column}{0.49\textwidth}
      % Earth measurement diagram
      \visualplaceholder[5cm]{Earth circle}
    \end{column}
  \end{columns}

  \note{
    This is one of the most elegant experiments in the history of science.
    With nothing but shadows, geometry, and a known distance, Eratosthenes
    measured the entire planet to within a few percent of the true value.
    Also mention the Sieve of Eratosthenes if time permits --- an algorithm
    for finding prime numbers that is still taught today.
  }
\end{frame}

% ----------------------------------------------------------
% Slide 21: Apollonius -- Conic Sections [EXTENDED]
% ----------------------------------------------------------
\begin{frame}{Apollonius --- Conic Sections}

  \begin{columns}[T]
    \begin{column}{0.48\textwidth}
      \begin{personbox}[Apollonius of Perga]{greekBlue}
        \small
        \textbf{Dates:} c.\ 262 -- c.\ 190 BCE\\
        \textbf{Origin:} Perga, Pamphylia (modern Turkey); worked in Alexandria\\
        \textbf{Contribution:} Wrote \emph{Conics} --- named the ellipse, parabola, and hyperbola
      \end{personbox}

      \vspace{0.4em}
      \begin{itemize}
        \item The definitive ancient work on conic sections
        \item \alert{These names are his}: ellipse, parabola, hyperbola
        \item Essential for Kepler's planetary orbits \textbf{1,800 years later}
      \end{itemize}

      \vspace{0.3em}
      \begin{insightbox}
        Apollonius named the curves you study in precalculus.
        Kepler would later show that planets travel in the ellipses Apollonius described.
      \end{insightbox}
    \end{column}

    \begin{column}{0.49\textwidth}
      % Conic sections diagram
      \visualplaceholder[5cm]{Draw cone outline}

      \vspace{0.3em}
      {\small\centering Slicing a cone at different angles\\produces all four curves.}
    \end{column}
  \end{columns}

  \note{
    If short on time, this slide can be cut [EXTENDED].
    Key point: pure mathematics done for its own sake in 200 BCE
    turned out to be exactly what Kepler needed in 1609.
    ``Math done for beauty often becomes math used for science.''
  }
\end{frame}

% ----------------------------------------------------------
% Slide 22: Diophantus -- Father of Algebra?
% ----------------------------------------------------------
\begin{frame}{Diophantus --- Father of Algebra?}

  \begin{columns}[T]
    \begin{column}{0.55\textwidth}
      \begin{personbox}[Diophantus of Alexandria]{greekBlue}
        \small
        \textbf{Dates:} fl.\ c.\ 250 CE\\
        \textbf{Origin:} Alexandria, Roman Egypt\\
        \textbf{Contribution:} Wrote \emph{Arithmetica} --- algebraic problems seeking
        integer or rational solutions; introduced proto-algebraic notation
      \end{personbox}

      \vspace{0.3em}
      \begin{itemize}
        \item \alert{Diophantine equations} (integer solutions) named after him
        \item One of the earliest uses of \textbf{symbolic abbreviations} for unknowns
        \item Fermat's ``Last Theorem'' was \emph{scribbled in the margin} of a copy of \emph{Arithmetica}
      \end{itemize}
    \end{column}

    \begin{column}{0.42\textwidth}
      \visualplaceholder[4.5cm]{Page from a manuscript of \emph{Arithmetica} (Wikimedia Commons)}

      \vspace{0.3em}
      % Fermat's margin note
      \visualplaceholder[5cm]{Diagram}
    \end{column}
  \end{columns}

  \note{
    \verifytag{Dates very uncertain (estimates range 200--284 CE).
    Qualify the claim that Diophantus ``introduced'' symbolic abbreviations ---
    earlier civilizations used some symbolic shorthand. His contribution was
    more systematic but not entirely unprecedented.}
    The Fermat connection makes this slide immediately relatable to students.
  }
\end{frame}

% ----------------------------------------------------------
% Slide 23: Hypatia of Alexandria
% ----------------------------------------------------------
\begin{frame}{Hypatia of Alexandria}

  \begin{columns}[T]
    \begin{column}{0.55\textwidth}
      \begin{personbox}[Hypatia of Alexandria (Hypatia h\={e} Alexandrin\={e})]{greekBlue}
        \small
        \textbf{Dates:} c.\ 350--370 -- 415 CE\\
        \textbf{Origin:} Alexandria, Roman Egypt\\
        \textbf{Contribution:} Mathematician, astronomer, philosopher.
        Commentaries on Diophantus and Apollonius.
        Head of the Neoplatonist school in Alexandria.
      \end{personbox}

      \vspace{0.4em}
      \begin{itemize}
        \item Wrote commentaries on Diophantus's \emph{Arithmetica} and Apollonius's \emph{Conics}
        \item One of the \textbf{last great scholars} of ancient Alexandria
        \item \alert{Murdered} by a Christian mob in 415 CE
        \item Her death is often seen as a symbol of the end of classical learning
      \end{itemize}
    \end{column}

    \begin{column}{0.42\textwidth}
      \visualplaceholder[4.5cm]{``Hypatia'' by Charles William Mitchell, 1885 (Laing Art Gallery, Newcastle)}

      \vspace{0.3em}
      % Timeline of Alexandria's decline
      \visualplaceholder[5cm]{Diagram}
    \end{column}
  \end{columns}

  \note{
    \verifytag{Birth year uncertain (c.\ 355 or 370 CE are common estimates).
    Whether she edited Arithmetica (wrote new commentary) or merely revised an
    existing commentary is debated.}
    ``She was one of the last great scholars of ancient Alexandria, and her murder
    is sometimes seen as a symbol of the end of classical learning.''
    Treat her story with respect --- she is a powerful example of a woman in mathematics
    in antiquity, but avoid reducing her to her death.

    PACING: After mentioning her murder, pause. Do not rush to the next slide.
    Let the silence speak.
  }
\end{frame}

% ----------------------------------------------------------
% Slide 24: Greek Legacy and Transition
% ----------------------------------------------------------
\begin{frame}{Greek Legacy and the Great Transition}

  \begin{columns}[T]
    \begin{column}{0.50\textwidth}
      \textbf{What the Greeks gave us:}
      \begin{itemize}
        \item \alert{Axiomatic reasoning} and \textbf{proof}
        \item The idea that math is about \emph{eternal truths}, not just practical calculation
        \item Geometry as a deductive system
        \item Foundations of number theory
      \end{itemize}

      \vspace{0.5em}
      \textbf{But then\ldots}
      \begin{itemize}
        \item Much Greek knowledge \textbf{lost} in the West
        \item \textbf{Preserved} in Byzantine and Islamic libraries
        \item Two very different journeys ahead\ldots
      \end{itemize}
    \end{column}

    \begin{column}{0.47\textwidth}
      % Knowledge flow diagram
      \visualplaceholder[5cm]{Greek source}

      \vspace{0.5em}
      {\small\textcolor{gray}{Greek texts survived thanks to\\Islamic and Byzantine scholars.}}
    \end{column}
  \end{columns}

  \vspace{0.3em}
  \begin{center}
    \textcolor{indianSaffron}{\emph{``Greek mathematical texts were about to take two very different journeys ---
    east to India and the Islamic world, and into near-oblivion in medieval Europe.''}}
  \end{center}

  \note{
    This is the transition slide. Emphasise that the Greek legacy was NOT
    a smooth handoff --- it was nearly lost. The Islamic world's role in
    preserving and extending Greek mathematics is CRUCIAL and will be covered
    in a later section.
    ``We are now at the end of Session 1. Next, we travel east.''
  }
\end{frame}
